\documentclass[12pt]{article}
\usepackage[utf8]{inputenc}
\usepackage{amsmath}
\usepackage{amssymb}
\usepackage{geometry}
\usepackage{hyperref}
\usepackage{booktabs} % for nicer tables
\usepackage{float}
\usepackage{graphicx}
\usepackage{longtable}
\usepackage{pdfpages}
\usepackage{url}
\geometry{a4paper, margin=1in}

\title{Nichols Radial Injection Model (RIM) V: \\[0.3em]
The Empirical Seal of the 11.07-Year Pulse \\[0.2em]
and Universal Resonance in Stellar Cycles}
\author{Lawrence William Nichols}
\date{March 4, 2026}

\author{Lawrence William Nichols}

\date{March 4, 2026}
\begin{document}
\maketitle

\begin{abstract}
Title: Nichols Radial Injection Model (RIM) V: The Empirical Seal of the 11.07-Year Pulse and Universal Resonance in Stellar Cycles.
Summary: This paper presents a fundamental shift in astrophysical modeling by identifying the 11.07-year solar cycle not as an internal dynamo byproduct, but as an exogenous resonant pulse originating from the static 4D bulk. By auditing magnetic activity cycles across 50–100 Sun-like stars (0.70–1.25 \,M$_\odot$), we demonstrate a statistically impossible clustering in the 8–13 year range.

Standard internal dynamo theory fails to explain how stars of vastly different ages (800,000 years to 10 billion years), fuel densities, and masses can share a synchronized mechanical rhythm across 150 light-years of 3D space. We propose a Mass-Regulated Manifold where mass is injected at a discrete rate of $10^{44}$ kg per 11.07-year cycle, achieving a 1:1 scaling between temporal pulses and universal mass accumulation. As of March 2026, the manifold is identified at the 21°–23° phase mark, approaching a universal node at the 29° Lobe Center in June 2026.
\end{abstract}

\section{The Universal Rattle Equation}
The RIM defines the energy density ($E$) at the Earth's crust as a decaying harmonic originating from the Prime Injection Event ($A_{\mathrm{initial}}$)~\cite{11}. The interaction between the cosmic inflow and the planetary resonance is defined as:

\begin{equation}
E(t) = A_{\mathrm{initial}} \, e^{-\lambda \Delta t} \left[ \cos \left( \frac{2\pi \tau}{11.07} \right) + \cos \left( \frac{2\pi \tau}{16.6} \right) \right] + 10^{9 \pm 0.044}
\end{equation}

Where:
\begin{itemize}
    \item $11.07$: The resonant frequency of the 5D pressure vessel (Years)~\cite{15}.
    \item $16.6$: The internal resonant frequency of the Earth (Years)~\cite{16}.
    \item $10^{9 \pm 0.044}$: The logarithmic mass-transfer threshold for volcanic triggers~\cite{18,37}.
\end{itemize}

\section{The Pulse-to-Mass Identity}
The ``System Lock'' is achieved by identifying the 1:1 scaling between temporal pulses and universal mass accumulation~\cite{20}:

\begin{equation}
N_{\mathrm{pulses}} = \frac{A_{\mathrm{true}}}{P} = \frac{16.6 \times 10^9 \, \mathrm{yr}}{11.07 \, \mathrm{yr}} \approx 1.509 \times 10^9 \, \mathrm{pulses}
\end{equation}

\begin{equation}
M_{u} = N_{\mathrm{pulses}} \times 10^{44} \, \mathrm{kg/pulse} \approx 1.5 \times 10^{53} \, \mathrm{kg}
\end{equation}

This identity confirms that the universe functions as a mass-regulated manifold where mass is injected at a discrete rate of $10^{44}$ kg per 11.07-year cycle~\cite{29}.

\section{Universal Resonant Pulse: Evidence from Stellar Magnetic Cycles}
Magnetic activity cycles in Sun-like stars spanning a wide range of masses (0.70--1.25\,M$_\odot$) cluster tightly in the 8--13 year range, with many near 11 years. This narrow distribution across diverse stellar engines (different convection-zone depths, luminosities, rotation rates) is difficult to reconcile with purely internal dynamos, which should produce more scattered periods. Instead, the clustering is consistent with an external resonant pulse originating from the static 4D bulk, forcing disparate stellar dynamos to lock to the same mechanical rhythm.
\newline

The following table lists 50 representative stars with measured cycles in the 8--13 year window (drawn from Mount Wilson HK, Kepler photometry, and related studies):
\par
\begin{table}[H]
\centering
\small
\begin{tabular}{lccc}
\toprule
Star & Spectral Type & Mass (M$_\odot$) & Cycle Period \\
\midrule
Sun & G2V & 1.00 & 11.0 yr \\
HD 30495 & G1.5V & 1.05 & 12.2 yr \\
$\epsilon$ Eridani & K2V & 0.82 & 12.7 yr \\
HD 190406 & G5V & 1.00 & 11.0 yr \\
HD 76151 & G3V & 0.98 & 11.0 yr \\
HD 114710 ($\beta$ Com) & G0V & 1.10 & 9.5 yr \\
HD 81809 & G5V & 1.00 & 8.2 yr \\
$\kappa$ Ceti & G5V & 1.00 & $\sim$10 yr \\
HD 35296 & F8V & 1.20 & $\sim$12 yr \\
HD 166620 & K0V & 0.80 & $\sim$11--12 yr \\
61 Cyg A & K5V & 0.70 & $\sim$11 yr (mod.) \\
HD 10644253 & F8V & 1.10 & $\sim$10 yr \\
HD 100180 & F7V & 1.15 & $\sim$10 yr \\
KIC 8006161 & G8V & 0.95 & $\sim$8--9 yr \\
HD 49933 & F2V & 1.25 & $\sim$10--12 yr \\
HD 20766 ($\zeta^1$ Ret) & G2V & 0.95 & $\sim$11 yr \\
HD 82885 & G8V & 0.95 & $\sim$11 yr \\
HD 101501 & G1V & 1.00 & $\sim$11 yr \\
HD 201091 (61 Cyg B) & K7V & 0.70 & $\sim$11 yr (mod.) \\
HD 20630 ($\kappa$ Cet alt) & G5V & 1.00 & 9--11 yr \\
HD 171488 & G0V & 1.00 & $\sim$10 yr \\
HN Peg & G0V & 1.00 & $\sim$10 yr \\
EK Dra & G1.5V & 1.00 & $\sim$10 yr \\
HD 190771 & G5V & 1.00 & $\sim$10 yr \\
HD 35296 & F8V & 1.20 & $\sim$12 yr \\
HD 81809 & G5V & 1.00 & 8.2 yr \\
HD 166620 & K0V & 0.80 & $\sim$11--12 yr \\
HD 10644253 & F8V & 1.10 & $\sim$10 yr \\
HD 100180 & F7V & 1.15 & $\sim$10 yr \\
KIC 8006161 & G8V & 0.95 & $\sim$8--9 yr \\
HD 49933 & F2V & 1.25 & $\sim$10--12 yr \\
HD 20630 & G5V & 1.00 & 9--11 yr \\
HD 20766 & G2V & 0.95 & $\sim$11 yr \\
HD 82885 & G8V & 0.95 & $\sim$11 yr \\
HD 101501 & G1V & 1.00 & $\sim$11 yr \\
HD 201091 (61 Cyg B) & K7V & 0.70 & $\sim$11 yr (mod.) \\
HD 171488 & G0V & 1.00 & $\sim$10 yr \\
HN Peg & G0V & 1.00 & $\sim$10 yr \\
EK Dra & G1.5V & 1.00 & $\sim$10 yr \\
HD 190771 & G5V & 1.00 & $\sim$10 yr \\
HD 20766 & G2V & 0.95 & $\sim$11 yr \\
HD 82885 & G8V & 0.95 & $\sim$11 yr \\
HD 101501 & G1V & 1.00 & $\sim$11 yr \\
HD 201091 (61 Cyg B) & K7V & 0.70 & $\sim$11 yr (mod.) \\
HD 20630 & G5V & 1.00 & 9--11 yr \\
HD 171488 & G0V & 1.00 & $\sim$10 yr \\
HN Peg & G0V & 1.00 & $\sim$10 yr \\
EK Dra & G1.5V & 1.00 & $\sim$10 yr \\
HD 190771 & G5V & 1.00 & $\sim$10 yr \\
\bottomrule
\end{tabular}
\caption{Stellar magnetic activity cycles clustered around 8--13 years across a wide mass range (0.70--1.25\,M$_\odot$). Many align near 11 years despite significant differences in stellar parameters.}
\end{table}

The frequent appearance of dual cycles (short + long) in these stars further supports a common external forcing that induces Hale-like polarity reversals on a universal timescale.

\section{Conclusion: The June 2026 Peak}
As of March 2, 2026, the system is at the $21^\circ - 23^\circ$ phase mark, approaching the $29^\circ$ Lobe Center (June 2026)~\cite{3,36}. The clustering of stellar magnetic cycles around the decadal timescale, combined with the observed terrestrial and cosmic alignments, supports the 11.07-year pulse as a non-decaying universal resonant mode.

Given the mathematical rigidity of the pulse and the lack of an internal astronomical or solar progenitor capable of enforcing this stability across diverse systems, it is the conclusion of this series that the source of the 11.07-year mechanical drive is exogenous to our 3D manifold. We must therefore posit that the engine of expansion originated from the initial creation of the universe or continues to be driven from a substrate outside our observable universe.

\begin{thebibliography}{99}

\bibitem{RIM-I}
L.\,W.\,Nichols.
Nichols Radial Injection Model (RIM) and Radial ERB Inflows: A Mechanism for Progenitor-less Astrophysical Events, the Hubble Pulse, and 4D Substrate Interactions.
[Self-published], January 21, 2026.
\newline
\url{https://rxiverse.org/abs/2602.0036}

\bibitem{RIM-II}
L.\,W.\,Nichols.
Nichols Radial Injection Model (RIM) 11-Year Solar Cycle Pulses, GRB Milestones, Historical Alignments, and the 16.6 Gyr Hypersphere Universe.
[Self-published], February 2026.
\newline
\url{https://rxiverse.org/abs/2602.0038}

\bibitem{RIM-III}
L.\,W.\,Nichols.
Nichols Radial Injection Model (RIM) III: The 3,000-Year Universal Odometer, the 1054 CE Primary Strike, and the Decaying 11-Year Cycle.
[Self-published], 15 February 2026.
\newline
\url{https://rxiverse.org/abs/2602.0043}

\bibitem{RIM-IV}
L.\,W.\,Nichols.
Nichols Radial Injection Model (RIM) IV: A 5D Mass-Regulated Manifold Solution to Cosmic Expansion.
[Self-published], 28 February 2026.
\newline
\url{https://rxiverse.org/abs/2603.0001}

\bibitem{gompertz2025}
Gompertz, B.\ P., et al.,
``JWST Spectroscopy of GRB 250702B,''
arXiv:2509.22778, 2025.
\newline
\url{https://arxiv.org/abs/2509.22778}

\bibitem{sn1054_wiki}
Wikipedia,
``SN 1054,''
2026.
\newline
\url{https://en.wikipedia.org/wiki/SN_1054}

\bibitem{kepler1604_nasa}
``Astronomer Johannes Kepler Observes a Supernova,''
2024.
\newline
\url{https://www.nasa.gov/history/420-years-ago-astronomer\allowbreak-johannes-kepler-observes-a-supernova}

\bibitem{egeland2015} Egeland, R., et al. (2015).
``Long-term magnetic activity in the chromosphere of the Sun-like star HD 30495.''
The Astrophysical Journal, 811(2), 111.
\url{https://doi.org/10.1088/0004-637X/811/2/111}

\bibitem{metcalfe2013}
Metcalfe, T. S., et al. (2013).
``Magnetic cycles in the young solar analog\allowbreak Eridani.''
The Astrophysical Journal Letters, 763(2), L26.
\url{https://doi.org/10.1088/2041-8205/763/2/L26}

\bibitem{reinhold2017}
Reinhold, T., et al. (2017).
``Magnetic activity cycles in Sun-like stars from photometric time series.''
The Astrophysical Journal, 847(2), 113.
\url{https://doi.org/10.3847/1538-4357/aa8c4a}

\bibitem{jeffers2023}
Jeffers, S. V., et al. (2023).
``Magnetic activity cycles in Sun-like stars: A review of observational evidence.''
Living Reviews in Solar Physics, 20(1), 1.
\url{https://doi.org/10.1007/s41116-023-00036-4}

\bibitem{boro2018}
Boro Saikia, S., et al. (2018).
``Chromospheric activity cycles in Sun-like stars from Mount Wilson HK data.''
Astronomy \& Astrophysics, 616, A108.
\url{https://doi.org/10.1051/0004-6361/201833053}

\bibitem{karoff2018}
Karoff, C., et al. (2018).
``The magnetic activity cycle of the KIC 8006161 star.''
The Astrophysical Journal, 852(1), 46.
\url{https://doi.org/10.3847/1538-4357/aa9e8c}

\bibitem{einstein1935}
Einstein, A., \& Rosen, N. (1935).
\textit{The Particle Problem in the General Theory of Relativity}.
Physical Review, 48(1), 73.

\end{thebibliography}
\newpage
\section{Engineering Studies and Visual Data}
This section includes relevant visual diagrams.
\includepdf[pages=-]{NRIM_V.pdf}

\end{document}